\section{The Model in Brief}
\label{sec:model}

We recall only what is needed to read the results. The full derivation is in Pollard (2026).

\begin{definition}[Vibe, tone, note, mesh]
A \emph{vibe} is the one kind of entity, a state in relation. Each vibe carries a finite \emph{tone}, with the working alphabet ternary, $t \in \{-1, 0, +1\}$, and richer spinor tones reserved for matter. The \emph{note} is the only primitive relation, $v$ noting $w$ meaning $v$ is present to $w$. The \emph{vibe mesh} is the whole population of vibes, their tones, and the note relation among them. There is no space or time outside the mesh.
\end{definition}

The mesh grows and updates by a local \emph{rule}, one \emph{beat} at a time, each vibe feeling only its neighbors. Four committed choices shape it, each with named alternatives held open:

\begin{itemize}
\item \textbf{Discreteness.} The base is finite. Real numbers appear only as emergent large-scale measurements, never in the substrate.
\item \textbf{Hyperbolic geometry.} The mesh branches outward exponentially rather than as a flat grid.
\item \textbf{Eternal expansion.} The mesh keeps growing without bound.
\item \textbf{No hidden state.} Everything is in the tones and notes. There is no separate bookkeeping behind them.
\end{itemize}

The causal order of the mesh, the relation of one vibe being in another's past, is a locally finite partial order, which is exactly a \emph{causal set} (Bombelli et al., 1987). This identification is what lets the geometry, dynamics, and matter of the mesh be studied with the established discrete-spacetime toolkit, and it is the bridge between Vibe Theory's ontology and the rest of this paper.

The guiding picture for the results is a ladder. At the bottom is the mesh of tones and notes. Above it, as large-scale patterns, are space, time, matter, force, gravity, and the correlations we call quantum. The paper climbs this ladder one rung at a time, and the final section shows how the rungs connect into a single structure.
